% !PRINTED 576
\weberpage{576}
%<<CONT
und folglich durch Vertauschung von $\xi$ mit $\xi'$:
\weq{11}{\theta_\xi(\tau_{\xi'}) = \theta_{\xi'}(\tau_\xi),}
und damit nach Bd. I, \S\ 169:

\begin{quote}
7. \gesperrt{Der Teilungskörper ist in bezug auf den Klassenkörper relativ Abelsch.}
\end{quote}

Die Zahlen $\xi$, nach dem Modul $\mathfrak{m}$ genommen, bilden bei der Multiplikation eine Gruppe, und die Relativgruppe des Teilungskörpers ist mit dieser Gruppe isomorph (oder wenigstens mit einem Teiler dieser Gruppe. Die Irreduzibilität von $T_{\mathfrak m}$ wird weiterhin bewiesen werden, wodurch dann die Gruppe von $\mathfrak{T}_{\mathfrak m}$ selbst festgestellt ist).

\begin{center}
\S\ 155. \textbf{Multiplikation der elliptischen Funktionen für einen}\\
\textbf{ungeraden Multiplikator.}
\end{center}

Für den Nachweis der Existenz des Teilungskörpers ist die Benutzung der \gesperrt{Weierstrass}schen Funktion und der daraus abgeleiteten $\tau$-Funktion sehr zweckmäßig, und es wäre am bequemsten, wenn man darauf auch die weitere Untersuchung dieses Körpers gründen könnte. Dafür aber ist die arithmetische Natur der Multiplikationsformeln noch nicht genügend bekannt, und es ist darum zurzeit noch notwendig, die komplexe Multiplikation und Teilung auch der \gesperrt{Jacobi}schen elliptischen Funktionen zu betrachten.

Hierbei wollen wir den Multiplikationsformeln, die wir in \S\ 57 abgeleitet haben, noch eine etwas andere Gestalt geben.

Wir betrachten zunächst durchweg den Multiplikator $m$ als \gesperrt{ungerade} Zahl. Wir setzen, wenn
\weq{1}{\varkappa = \frac{\vartheta_{10}}{\vartheta_{00}}}
den Modul der elliptischen Funktionen bedeutet, nach \gesperrt{Kronecker}\textsuperscript{1)}\begingroup\renewcommand{\thefootnote}{1)}\footnotetext[1]{Zur Theorie der elliptischen Funktionen. Sitzungsbericht der Berliner Akademie vom 29. Juli 1886.}\endgroup:
\weq{2}{\lambda = 4\left(\varkappa + \frac{1}{\varkappa}\right) = 4\,\frac{f(\omega)^4}{f_2(\omega)^4} + 4\,\frac{f_2(\omega)^4}{f(\omega)^4} \quad [\S\ 54, (3)],}
und entwickeln diesen Ausdruck nach steigenden Potenzen von
\[
q = e^{\pi i \omega}.
\]
